% Text overlay only; model artwork is retained without modification.
\begin{tikzpicture}[x=1mm,y=1mm,text=BookInk]
\path[use as bounding box] (0,0) rectangle (157.4,70.43125333333333);
\node[inner sep=0pt,outer sep=0pt,align=center,anchor=west,font=\FigureNoteFont\fontsize{7.5}{9.5}\selectfont,text width=53.516000000000005mm] at (7.8700,66.5575) {(a) 显式状态模型};
\node[inner sep=0pt,outer sep=0pt,align=center,anchor=center,font=\figurefont\fontsize{8}{10}\selectfont,text width=27.544999999999998mm] at (18.5732,51.1331) {状态$s_t$、动作$a_t$};
\node[inner sep=0pt,outer sep=0pt,align=center,anchor=center,font=\figurefont\fontsize{8}{10}\selectfont,text width=27.544999999999998mm] at (59.4972,51.1331) {条件动力学$\widehat T_{\bm\phi}$};
\node[inner sep=0pt,outer sep=0pt,align=center,anchor=center,font=\figurefont\fontsize{8}{10}\selectfont,text width=27.544999999999998mm] at (98.5324,51.1331) {后继分布\\$\widehat T_{\bm\phi}(\cdot\mid s_t,a_t)$};
\node[inner sep=0pt,outer sep=0pt,align=center,anchor=center,font=\figurefont\fontsize{8}{10}\selectfont,text width=26.758000000000003mm] at (139.4564,51.1331) {奖励$\widehat r_t$\\继续概率$\widehat c_t$};
\node[inner sep=0pt,outer sep=0pt,align=center,anchor=west,font=\FigureNoteFont\fontsize{7.5}{9.5}\selectfont,text width=78.7mm] at (7.8700,38.5963) {(b) 从观测历史形成潜在状态};
\node[inner sep=0pt,outer sep=0pt,align=center,anchor=center,font=\figurefont\fontsize{8}{10}\selectfont,text width=27.544999999999998mm] at (18.5732,26.2004) {观测与动作历史\\$o_{\leq t},a_{<t}$};
\node[inner sep=0pt,outer sep=0pt,align=center,anchor=center,font=\figurefont\fontsize{8}{10}\selectfont,text width=27.544999999999998mm] at (59.4972,26.2004) {编码或推断$\bm z_t$};
\node[inner sep=0pt,outer sep=0pt,align=center,anchor=center,font=\figurefont\fontsize{8}{10}\selectfont,text width=27.544999999999998mm] at (98.5324,26.2004) {潜在转移\\$(\bm z_t,a_t)\to\bm z_{t+1}$};
\node[inner sep=0pt,outer sep=0pt,align=center,anchor=center,font=\figurefont\fontsize{8}{10}\selectfont,text width=26.758000000000003mm] at (139.4564,26.2004) {奖励与继续预测\\任务预测头};
\node[inner sep=0pt,outer sep=0pt,align=center,anchor=center,font=\figurefont\fontsize{8}{10}\selectfont,text width=27.544999999999998mm] at (59.4972,8.4518) {观测解码$\widehat o_t$};
\node[inner sep=0pt,outer sep=0pt,align=center,anchor=center,font=\FigureNoteFont\fontsize{7.5}{9.5}\selectfont,text width=15.740000000000002mm] at (78.7000,29.5811) {动作$a_t$};
\node[inner sep=0pt,outer sep=0pt,align=center,anchor=center,font=\FigureNoteFont\fontsize{7.5}{9.5}\selectfont,text width=64.53399999999999mm] at (116.4760,8.4518) {重建帮助保留观测；控制还需动作条件预测与充分表示。};
\end{tikzpicture}
